% Independent audit of the scalar-predecessor dependency boundary.
\section{Predecessor in the fragment without witnessed introduction}
\label{sec:predecessor-fragment}

This appendix strengthens the dependency statement of the inherited scalar
predecessor construction.  Its entire proof can be conducted using equations,
the two iterator computations, ordinary equality reasoning, the progressive
closure rule for literal equations, and closure specialization.  In particular,
it needs no instance of witnessed introduction.  This matters because this
fragment is closed under arbitrary numerical substitution, even when the full
calculus restricts its existential witnesses to variables and zero.

Write $C(n,a,b)=I(n,a;u.b)$, $Z(n)=C(n,1,0)$, and
$a+b=I(b,a;u.Su)$.  Let $\mathcal E$ denote the fragment just described.
An equation-case inference belongs to $\mathcal E$: from $E(0)$ and a proof
of $E(Sz)$ for fresh $z$, add the unused assumption $E(z)$ in the latter
proof, apply progressive closure, and specialize it using $E(0)$.
This inference concerns an equation $E$; it does not presuppose a
disjunction $t=0\lor\exists z\,t=Sz$.

The arithmetic equations for a commutative monoid follow by equation
induction.  Equation cases give the two unconditional identities
\[
 Z(v)+Z(Z(v))=1,\qquad Z(Z(Z(v)))=Z(v).
\]
For $N(n)=I(n,0;u.Z(u))$, equation induction consequently gives
\[
 N(n)+Z(N(n))=1,\qquad Z(Z(N(n)))=N(n).
\]
Thus every later use of the statement that $N(n)$ is Boolean can be
replaced by explicit equations.  In particular $N(n+2)=N(n)$, followed
by $N(2n)=0$, $N(2n+1)=1$, and $N(3n)=N(n)$, uses only equation induction.

The potentially delicate point is the ramp invariant.  Put
\[
 j(u)=Z(u)+2N(u),\qquad M(n)=I(n,0;u.j(u)),\qquad
 G(n)=I(n,0;u.u+M(u)+1).
\]
Finite computations give $j(0)=1$, $j(1)=2$, and $j(2)=0$.
The following two unconditional equations are established by equation
cases on the fresh variable $v$:
\begin{align*}
 M(g+Z(v)+1)&=C(v,j(j(M(g))),j(M(g))),\\
 C(v,j(j(Z(v))),j(Z(v)))&=Z(Z(v)).
\end{align*}
For the first, the zero branch is the double successor computation for
$M$, and the successor branch is its single successor computation.
For the second, the zero branch computes $j(j(1))=0$ and the successor
branch computes $j(0)=1$.

Now suppose the single induction equation is $M(G(n))=N(n)$.
Write $g=G(n)$ and $p=N(n)$.  Substitute $v=Z(p)$ into the two
displayed equations and use $Z(Z(p))=p$.  Ordinary equality reasoning
gives $M(g+p+1)=Z(p)$.  The update for $G$ and the current induction
equation identify $G(Sn)$ with $g+p+1$, while $N(Sn)=Z(p)$.
This proves $M(G(n))=N(n)$ by equation induction with no conditional
Boolean splitting.

The invariant gives $G(Sn)=G(n)+N(n)+1$ and
$G(n+2)=G(n)+3$, the latter using $N(n)+Z(N(n))=1$.
Consequently equation induction proves $G(2n)=3n$.  The successor
calculation for $G(3n)=3G(n)+N(n)$ reduces, with $p=N(n)$, to
\[
 2p+4=3p+3+Z(p).
\]
This is the monoid equation obtained from $p+Z(p)=1$; it requires no
disjunctive Boolean range assertion.

Finally put $H(u)=C(N(u),G(u),u)$.  The remaining apparent case split
in $H(3u)=3H(u)$ also has a purely equational replacement.  Equation
cases prove, for arbitrary parameters,
\[
 C(e,a+b,c+d)=C(e,a,c)+C(e,b,d),\qquad C(e,e,0)=0.
\]
Using $N(3u)=N(u)$ and $G(3u)=3G(u)+N(u)$, these yield
\begin{align*}
 H(3u)
 &=C(N(u),3G(u)+N(u),3u)\\
 &=3C(N(u),G(u),u)+C(N(u),N(u),0)=3H(u).
\end{align*}
The equations $H(2u)=3u$ and $H(1)=1$ are immediate.

From this point the inherited comparator and saturation construction
uses only equation induction and equality: the semigroup and rotation
identities for iteration; $I(y,3a;H)=3I(y,a;H)$ and
$I(y,1;H)=1$; the comparator
$L(x,y)=N(I(y,I(x,1;u.2u);H))$; and the saturated iterator
$W(t,m)=I(t,0;u.C(L(Su,m),u,Su))$.
Its only subsequent numerical case, the identity
$C(e,Sa,Sb)=S(C(e,a,b))$, is itself an equation case.
The final finite scalar term
\[
 P(m)=I(m,0;u.C(L(SSu,m),u,Su))
\]
therefore satisfies
\[
 \mathcal E\vdash P(0)=0,\qquad
 \mathcal E\vdash P(Sm)=m.
\]
The whole derivation, not merely its conclusion, is stable under numerical
substitution.  Applying the already derived term $P$ to $Sx=Sy$ gives
$x=y$ in the same fragment.
